\subsection{Finite geometric systems and their limits}
\label{subsec:eight-finite-constructions}

The preceding constructions have different meanings of finiteness.
A bounded set of charges is combinatorial. A moduli stack of finite
type is geometric. A finite-dimensional coefficient space is a
property of a specified cohomology theory. None of these conditions
alone supplies the other two. For example,
\(H^*(B\mathbb G_m;\mathbb Q)=\mathbb Q[c_1]\) is unbounded even
though \(B\mathbb G_m\) is an Artin stack of finite type.

We organize the compact comparison by finite operation diagrams and
then by inverse systems. A finite stage must contain the input and
output spaces of the operations being compared. It need not be
closed under every extension or bracket. When a calculation leaves
the retained range, its larger output stage must be specified.
This permits finite calculations without asserting that a bounded
root fragment is a Lie algebra or that a bounded nonzero charge set
is closed under arbitrary addition.

\subsubsection*{Charges and target windows}

Fix the stability and support data of
Appendix~\ref{app:stability-support}. Let
\(F_{\sigma,S}(R)\subset\Gamma_X^{\mathrm{form}}\) be a supplied
finite set of charges with positive height bounded by \(R\), and
let \(\Gamma_R^{HN}\subset F_{\sigma,S}(R)\) be the charges of
the geometric strata to be used. The product-stability construction
has central charge
\[
 L_F(n)=Z_S\bigl(Rp_{S,*}(F\otimes p_E^*H^n)\bigr),
\]
where \(H\) has degree one on \(E\) and the K3 stability datum is
rational \cite{LiuProductStability}. Its numerical formula does not
prove boundedness of the chosen moduli problem. The relevant class
bounds, bounded families and finite-type substacks must be supplied
at their actual support and stability scope.

Let \(\Pi_X\) and \(B\) be the Gram map and normal-ordering
cocycle of Section~\ref{subsec:d6-d2-d0-dictionary-appendixC}.
The cocycle equation is
\[
 B(c,c')+B(c+c',c'')=B(c',c'')+B(c,c'+c'').
\]
It makes the central lift associative on its declared charge monoid.
Restrictions to finite sets retain only the products whose outputs
are present.

\begin{definition}[Finite target test window]
\label{def:finite-target-test-window}
For \(c\in\Gamma_R^{HN}\), set
\[
 \mathcal T_R(c)=
 \left\{\sum_{i<j}B(c_i,c_j):
   c=\sum_i c_i,\ c_i\in F_{\sigma,S}(R),\
   \sum_i h_S(c_i)\le R\right\}.
\]
Assume these translate sets are finite, and define
\[
 \begin{split}
 \widehat\Gamma_R&=
    \{(c,T):c\in\Gamma_R^{HN},\ T\in\mathcal T_R(c)\},\\
 \Gamma_R^{\mathrm{test}}&=
    \{\Pi_X(c)+T:(c,T)\in\widehat\Gamma_R\}.
 \end{split}
\]
This is a finite target-degree set. Compatible inclusions of the
charge and translate sets give inclusions of the target sets as
\(R\) increases.
\end{definition}

The image formula constructs a set of Gram degrees after its inputs
have been supplied. It does not construct sheaves representing those
degrees, nor does it make the test window exhaustive.

\begin{definition}[Exhaustion of a specified target set]
\label{def:relation-closed-target-finite-r-exhaustion}
Let \(\mathcal D_{\mathrm{rel}}\) be a specified set of degrees
needed by the target relations. A finite-stage realization of this
set assigns to each \(\gamma\in\mathcal D_{\mathrm{rel}}\) a
finite height \(R(\gamma)\), a charge
\(c_\gamma\in\Gamma_{R(\gamma)}^{HN}\), and a translate
\(T_\gamma\in\mathcal T_{R(\gamma)}(c_\gamma)\), such that
\[
 \Pi_X(c_\gamma)+T_\gamma=\gamma.
\]
The assignments are compatible with the stage inclusions. If the
target set has a finite prefix and symbolic tail families
\(\gamma_N\), the tail requires formulas
\(R_N,c_N,T_N\) valid for every allowed \(N\), with the same
image identity. All intermediate degrees of each asserted relation
must be included in an appropriate finite enlargement.
\end{definition}

\begin{proposition}[Coverage of target degrees]
\label{prop:relation-closed-target-finite-r-exhaustion-criterion}
Under this datum, every degree of \(\mathcal D_{\mathrm{rel}}\)
belongs to \(\Gamma_R^{\mathrm{test}}\) for some finite \(R\).
\end{proposition}
\begin{proof}
The pair \((c_\gamma,T_\gamma)\) belongs to
\(\widehat\Gamma_{R(\gamma)}\) and has image \(\gamma\).
The same argument applies uniformly to each symbolic family.
\end{proof}

A relation cascade, including a family of the form
\(A_N\Rightarrow B_N\Rightarrow A_{N+2}\), can describe target
degrees that must be retained. It supplies neither the geometric
charges \(c_N\) nor their normal-ordering translates. Conversely,
a finite list of represented degrees does not prove coverage of an
infinite cascade. The target calculations and the geometric
exhaustion are distinct parts of recognition.

\subsubsection*{Geometric stages and transitions}

A geometric stage consists of the chosen finite-type object stacks,
extension stacks and flag stacks, with the support and stability
conditions needed for their maps. The full CY3 deformation problem
has its shifted symplectic structure on its stated derived stack
\cite{PTVV2013}. Restricting or reducing that problem requires a
comparison preserving the relevant structure; a finite-type
substack is not automatically a two-term critical chart.
The reduced constructible coefficients and orientations are those
constructed on the actual reduced critical problem
\cite{BBDJS2015,JoyceUpmeier}.

A universal complex may live on a scalar gerbe. Its weight-zero
endomorphism complex and natural maps can descend without descent
of the weight-one universal object. Scalar rigidification likewise
does not imply finite residual inertia: direct sums can retain
positive-dimensional automorphism groups. Both facts affect the
coefficient theory and its quotient maps.

\begin{definition}[Finite geometric transition datum]
\label{def:finite-moduli-transition-geometry}
For \(R'\le R\), specify object, extension and flag morphisms
\[
 t^a_{RR'}:\mathfrak M^a_R\longrightarrow\mathfrak M^a_{R'},
 \qquad a\in\{\mathrm{obj},\mathrm{ext},\mathrm{fl}\},
\]
and closed retained substacks
\(\mathfrak Z^a_R\subset\mathfrak M^a_R\) carried into
\(\mathfrak Z^a_{R'}\). Require the restricted morphisms to be
proper; a closed immersion is one allowed special case.
For every coefficient operation, include the appropriate base-change
square and its comparison morphism. Include the maps of deformation
complexes, cosections, constructible coefficients and root pairs,
with their overlap and composition coherences.
\end{definition}

\begin{proposition}[Restricted geometric transitions]
\label{prop:finite-moduli-transition-geometry-criterion}
These data give proper transition maps on the retained substacks.
When their coefficient and base-change comparisons are equivalences,
they transport the specified coefficient operations and orientation
pairs compatibly with composition.
\end{proposition}
\begin{proof}
The retained-locus condition restricts each supplied morphism to the
chosen source and target. Properness is a hypothesis on those maps.
The coefficient assertion follows by composing their specified
pullback or pushforward comparisons with the base-change
isomorphisms. Their coherence gives equality of the composite
transports.
\end{proof}

A closed graph alone does not make a morphism a closed immersion.
Nor do inclusions of bounded HN substacks naturally give morphisms
in the opposite direction. The direction of a transition and the
variance of its coefficient functor are part of the datum. A direct
system obtained from open inclusions is not the inverse system above
without an additional construction.

\subsubsection*{The eight kinds of structure}

The following structures specify the operations that a finite
comparison must retain.
\begin{enumerate}
\item \emph{Charge and support.} The geometric classes, Gram map,
  normal-ordering cocycle and finite operation degrees are those
  above. The elliptic support degree is defined by
  \((p_E)_*\operatorname{cyc}_1(F)=b^{\mathrm{geom}}[E]\).
  It is additive on the relevant exact sequences. It is not the
  third Borcherds coordinate \(m^{\mathrm{Bch}}\) without a proved
  comparison. A positive effective elliptic degree projects onto
  all of \(E\); at degree zero the image of the curve support
  in \(E\) is a finite set.
\item \emph{Moduli and coefficients.} Specify the derived object
  and extension stacks, full flag correspondences, reduced
  deformation theory, critical coefficient system and supports.
  A proper target map can be obtained from a closed relative Quot
  locus in the sheaf recipient when its boundedness and closedness
  hypotheses hold. It does not imply properness of the source map,
  and the argument does not apply to an arbitrary tilted-heart
  extension problem without another construction.
\item \emph{Hybrid operations.} The local and wrapped sectors of
  Section~\ref{subsec:hybrid-ran-carrier} use their actual
  configuration spaces, wrapped prequotients and anchors. The
  \(LL,LW,WL,WW\) correspondences retain the source and target
  anchor data before quotient. Their three-input comparisons include
  all eight words in \(\{L,W\}^3\), and the full higher-flag
  coherences are further structure. The wrapped product of quotient
  stacks uses a common \(BE\) fibre, not two independently
  forgotten translation torsors. The quotient-after-correspondence
  pseudofunctor and protected integration maps must satisfy the
  relevant base-change, projection and Thom--Sebastiani identities.
\item \emph{Orientation and transport.} Specify root pairs on the
  actual quotient, cross-determinant factors on extensions and
  coherent flag comparisons. On a finite-stabilizer locus, a root
  obstruction and the character of a chosen root are separate
  conditions. Weyl lifts act on this same datum and preserve its
  square maps, critical coefficients and transition comparisons.
\item \emph{Hall and primitive operations.} Specify product,
  coproduct, unit, counit and the tensor in which the bialgebra
  equation is asserted. The ordinary Koszul and meromorphic
  braided tensors give different equations. A product alone gives
  neither cooperation nor a primitive projection. In a connected
  cocommutative ordinary realization over characteristic zero,
  the Eulerian projection has the hypotheses of
  Section~\ref{subsec:primitive-recognition}. The invariant Lie
  form, its radical and its generated Hopf ideal are distinct from
  any separately supplied full Hopf pairing.
\item \emph{Pfaffians.} Supply the finite skew complexes, their
  shifts and square maps, the Pfaffian lines and sections, and the
  geometric wall charts. The actual parameter and automorphic-line
  comparisons fix the target and its normalization. Positive-height
  product convergence and analytic continuation are separate from
  geometric effectivity of the family.
\item \emph{Associative coefficients.} Supply the counital,
  coaugmented source coalgebra with its conilpotence or the precise
  replacement filtration, and the associative cobar recipient of
  Section~\ref{subsec:koszul-source}. The finite comparison maps
  preserve its differential, support, tensor and operations.
  Lie--commutative chiral duality is not a substitute for this
  associative comparison.
\item \emph{Recognition.} Supply geometric simple and Cartan
  representatives, relation maps, kernel equalities, generation,
  both root parities and PBW comparisons. A finite calculation must
  retain all intermediate outputs. Compatible comparisons in every
  bounded weighted word space give the all-level algebraic
  assertion; a topological conclusion also requires continuous
  inverse maps on a specified completion.
\end{enumerate}

For wrapped sheaves, a determinant anchor can record translation
when its Euler characteristic is nonzero. At Euler characteristic
zero that information can vanish. Jordan--H\"older data may then
separate additional objects but do not automatically determine an
elliptic translation orbit. Additional anchors must be chosen and
their separating property proved. Similarly, a change of origin can
change a line's Picard class without determining an equivariant
stabilizer character. The constructions and examples of
Section~\ref{subsec:hybrid-ran-carrier} keep these distinctions.

On the formal nilpotent support of the local deformation problem,
the cyclic local model and
its full flag extensions provide genuine local operations. Their
comparison with a global reduced compact coefficient functor, its
proper correspondence maps and quotient descent remains a separate
geometric construction. The finite organization above preserves
this positive local result without treating a list of required
objects as proof of global existence.

\subsubsection*{Morphisms of finite comparison data}

\begin{definition}[Finite-stage Dirac--Igusa datum and morphism]
\label{def:finite-stage-dirac-igusa-morphism}
In a fixed coefficient setting, a stage is a tuple
\[
 \begin{split}
 \mathfrak D^{\mathrm{DI}}_{X,R}=
 (&\mathcal A_{X,R}^{E_3},F_{X,R}^{\mathrm{hyb}},
   \Gamma_{X,R},\Pi_X,\Phi_R,o_R,H_R,P_R^\Pi,\\
  &C_{X,R},\Theta_{\mathrm{Kos},R},
   \mathscr L_{\mathrm{Pf},R},\operatorname{pf}_{X,R},
   \epsilon_{o,R},\mathrm{Rec}_R),
 \end{split}
\]
whose components carry the specified structures above. In particular,
an \(E_3\)-source is actual additional operadic data when that
component is required; it is not inferred from its scalar series.
A morphism consists of the following compatible maps.
\begin{enumerate}
\item[(M1)] Maps of charge and support diagrams preserving
  \(\Pi_X\), the cocycle \(B\) and the retained products.
\item[(M2)] Geometric and coefficient transition maps as in
  Definition~\ref{def:finite-moduli-transition-geometry}, with the
  required proper support and variance.
\item[(M3)] Maps of the hybrid carriers and all specified
  coloured-flag operations, compatible with the quotient
  pseudofunctor and protected integration.
\item[(M4)] Isomorphisms of the pulled-back orientation pairs,
  compatible Weyl lifts and wall charts, preserving their square
  maps, normal coordinates and section comparisons.
\item[(M5)] Maps of the chosen Hall data preserving product,
  coproduct, units and counits in their declared tensor. Preserve
  primitives, the invariant Lie form and its radical; require a
  further Hopf-pairing comparison only where that pairing is
  independently supplied.
\item[(M6)] Isomorphisms of pulled-back Pfaffian lines carrying
  the specified sections to one another and commuting with their
  square and automorphic-line comparisons.
\item[(M7)] Coalgebra maps preserving the differential, counit,
  coaugmentation, filtration and the associative comparison
  \(\Theta_{\mathrm{Kos}}\).
\item[(M8)] Maps of the recognition diagrams preserving the
  representatives, relations, kernels, parity pieces, invariant
  forms, radicals and PBW filtrations, with all required
  intermediate operation spaces.
\end{enumerate}
Composition is componentwise, with the coherence required by the
chosen coefficient category. Fix strict models when an ordinary
category \(\mathsf C_X\) is used below. Alternatively retain an
enhanced category and all its mapping spaces. A finite-stage
isomorphism is an invertible morphism of these structured data,
not merely an isomorphism of their underlying vector spaces.
\end{definition}

The preservation results in the preceding subsections apply to
these maps at their individual scopes: orientation pairs and
critical coefficients, Hall products and coproducts, primitive
subspaces, invariant-form radicals, PBW filtrations and parity.
They do not follow from a scalar equality. In particular, maps
between nonorthogonal matrix quotients must be checked before a
radical or pairing is transported.

\subsubsection*{Pro-objects and actual inverse limits}

\begin{definition}[The Dirac--Igusa pro-object]
\label{def:dirac-igusa-pro-isomorphism}
A compatible cofiltered diagram
\(R\mapsto\mathfrak D^{\mathrm{DI}}_{X,R}\) in \(\mathsf C_X\)
defines the pro-object
\(\mathfrak D_X^{\mathrm{DI}}=\{\mathfrak D^{\mathrm{DI}}_{X,R}\}_R\).
A pro-isomorphism is an invertible morphism in
\(\operatorname{Pro}(\mathsf C_X)\).
No componentwise inverse limit is implicit in this notation.
\end{definition}

\begin{proposition}[Mapping formula and cofinal invariance]
\label{prop:dirac-igusa-pro-object-universal-property}
For two such systems \(\mathfrak E=(\mathfrak E_S)_S\) and
\(\mathfrak D=(\mathfrak D_R)_R\),
\[
 \operatorname{Hom}_{\operatorname{Pro}(\mathsf C_X)}
       (\mathfrak E,\mathfrak D)
 =\varprojlim_R\varinjlim_S
       \operatorname{Hom}_{\mathsf C_X}(\mathfrak E_S,\mathfrak D_R).
\]
For a constant source \(T\), this reduces to
\[
 \operatorname{Hom}_{\operatorname{Pro}(\mathsf C_X)}(T,\mathfrak D)
 =\varprojlim_R\operatorname{Hom}_{\mathsf C_X}(T,\mathfrak D_R).
\]
Cofinal restriction gives an isomorphic pro-object. Compatible
finite-stage isomorphisms on common cofinal systems are a sufficient
construction of a pro-isomorphism.
\end{proposition}
\begin{proof}
The first formula defines the morphisms of the pro-category. A
constant source removes the inner colimit. Cofinal restriction
preserves these mapping functors, hence the represented pro-object.
Compatible finite-stage inverses give inverse elements in the two
mapping sets, proving the last assertion. In an enhanced
pro-category, the corresponding statement uses mapping spaces and
homotopy limits and colimits.
\end{proof}

The sufficient condition is not necessary in a general pro-category.
For example, the vector-space tower
\[
 \mathbb C\xleftarrow{0}\mathbb C\xleftarrow{0}\mathbb C
       \xleftarrow{0}\cdots
\]
is pro-zero: for every target stage, the next transition kills the
identity class in the defining colimit. Every stage is nevertheless
nonzero, so no cofinal set of stages is levelwise isomorphic to the
zero tower. Replacing the zero transitions by identities gives the
constant nonzero pro-object. Thus the transition maps are decisive.

When an actual inverse limit exists, it belongs to a specified
complete category. For modules its projection topology has basic
neighbourhoods given by kernels of the stage projections. For line
or root-pair components the appropriate construction is a limit of
Picard groupoids or stacks, with its descent data. Neither existence
of a pro-object nor Mittag--Leffler vanishing for some cohomology
groups proves geometric effectivity of all its components.

\begin{definition}[Degreewise Mittag--Leffler condition]
\label{def:moduli-cohomology-ml-datum}
Choose a countable cofinal sequence and coefficient complexes
\(K_{R,a}\) on the retained object, extension, wrapped or flag
stacks. Let
\[
 C^q_{R,a}=H_c^q(\mathfrak M_{R,a},K_{R,a})
\]
when compactly supported cohomology is the selected theory, and use
the specified dualizing-complex model in a Borel--Moore theory.
The actual six-functor maps define a tower only after their variance
and identity and composition laws have been fixed.
It is Mittag--Leffler in degree \(q\) if, for every \(R_0\), the
images of \(C^q_{R,a}\to C^q_{R_0,a}\) are constant for all
sufficiently large \(R\).
\end{definition}

\begin{proposition}[Degreewise exactness of a countable limit]
\label{prop:moduli-cohomology-lim1-vanishing-criterion}
For a countable tower of modules satisfying the preceding condition,
\(R^1\varprojlim_R C^q_{R,a}=0\).
For a tower of complexes \(C_R\), there is a short exact sequence
\[
 0\longrightarrow\varprojlim{}^1 H^{q-1}(C_R)
 \longrightarrow H^q(R\varprojlim C_R)
 \longrightarrow\varprojlim H^q(C_R)\longrightarrow0.
\]
In particular, degreewise Mittag--Leffler cohomology gives the
right-hand identification.
\end{proposition}
\begin{proof}
The derived limit of a countable module tower is represented by the
two-term product complex with map
\[
 1-T:\prod_n A_n\longrightarrow\prod_n A_n,
 \qquad (1-T)(x)_n=x_n-u_nx_{n+1}.
\]
For a surjective tower this map is surjective by successive lifting.
For a Mittag--Leffler tower, its stabilized images form a
surjective subtower. The quotient is pro-zero. On a pro-zero tower,
\(\sum_{j\ge0}T^j\) is finite in each fixed output coordinate
and inverts \(1-T\). The exact sequence of these towers therefore
gives surjectivity of \(1-T\) on the original tower, proving the
vanishing of its cokernel \(\varprojlim{}^1\).
Applying the same product-cone model to a tower of complexes gives
the displayed cohomology sequence.
\end{proof}

Every tower of finite-dimensional vector spaces over a fixed field
is Mittag--Leffler: its images in a fixed stage form a descending
chain of subspaces. The analogous assertion for finite-rank
integral modules is false, as the tower
\(\mathbb Z\xleftarrow{2}\mathbb Z\xleftarrow{2}\cdots\)
shows. The images \(2^n\mathbb Z\) never stabilize.
Finite type of a stack, a proper transition map or an Euler
characteristic does not establish the finite-dimensional or
Mittag--Leffler hypotheses for a chosen coefficient theory.

\subsubsection*{Finite and limiting comparison statements}

\begin{theorem}[Formation of the finite comparison system]
\label{thm:finite-stage-dirac-igusa-pro-object}
Suppose the eight kinds of structure above are constructed, with
morphisms \textup{(M1)}--\textup{(M8)} satisfying their identity and
composition laws on a cofiltered stage system. They define a
Dirac--Igusa pro-object. The primitive, pairing-kernel, PBW and
coefficient cohomology limits agree with the cohomology of their
derived limits whenever each of those towers satisfies the
Mittag--Leffler hypotheses just stated. No geometric effectivity or
analytic continuation is implied by this conclusion.
\end{theorem}
\begin{proof}
The component maps and their coherence define the diagram in
\(\mathsf C_X\), hence its pro-object. Apply the preceding
proposition to each of the indicated coefficient towers. The last
sentence distinguishes these limits from the further geometric and
analytic realizations.
\end{proof}

\begin{theorem}[Finite Pfaffian comparison]
\label{thm:pfaffian-dirac-finite-stage}
At a finite stage carrying the geometric skew and orientation data,
the Pfaffian product and local reflection identities hold under
their specified line comparisons. A primitive recognition or
associative cobar conclusion holds only with the corresponding
finite source maps and relations. If these comparisons and their
inverses commute with every stage transition, they give the stated
comparison of pro-objects and, when defined, of their compatible
limits.
\end{theorem}
\begin{proof}
The finite skew identity is
\(\operatorname{pf}(D)^2=\det D\), and orthogonal block sums
multiply Pfaffians. Section-compatible wall lifts give the local
sign of Proposition~\ref{prop:rank-one-crossing-o2-sign}.
The primitive and coefficient conclusions are the separate
recognition and associative comparison theorems of the preceding
subsections. Compatible inverse maps induce a pro-isomorphism by
Proposition~\ref{prop:dirac-igusa-pro-object-universal-property};
applying an available limit functor preserves that equivalence.
\end{proof}

For \(K3\times E\), the finite moduli, reduced coefficients,
orientation transport and compact comparison maps remain the
geometric inputs. The local cyclic charts, formal charge formulas,
target relations and finite skew identities provide concrete parts
of the construction. A global realization requires their compatible
geometric extensions, rather than only their scalar shadows.
